\section{单卡预训练与规模实验}
\label{sec:pretrain}

\subsection{训练程序和资源配置}

预训练均从随机权重开始，用 RTX 4080 SUPER 在 Ubuntu 22.04、Python 3.11.16 和 PyTorch 2.14.1+cu126 环境运行。正式训练开启 BF16 自动混合精度，保留 FP32 主参数。优化使用 AdamW~\cite{loshchilov2019decoupled}，$\beta_1=0.9$、$\beta_2=0.95$、$\epsilon=10^{-8}$，二维及以上权重应用权重衰减。梯度全局范数限制为 1.0，学习率先预热再按余弦变化。唐诗等预算组用微批 32、累积 4，完整优化步的目标数为 $32\times4\times256=32{,}768$。计算每个累积窗口的有效目标总数以后再归一化梯度，可避免短窗和末尾批次改变优化权重。

训练程序在检查点保存模型参数、优化器、调度器、随机数状态、数据游标及数据与代码指纹。续训从游标定位到同一轮尚未处理的窗口。实际 14M 基线在第 87 步落盘，然后恢复到第 870 步。唐宋 92M 的长程训练也从第 1,974 步继续，验证日志和最优权重可以对应到累计有效目标数。这使一次长训练可以分多次运行，统计仍使用统一口径。

\subsection{固定训练预算下扩大模型}

14M、42M 和 92M 使用同一唐诗语料、同一字表、10 轮及约 2,849 万有效目标。14M 的验证 NLL 从 9.1361 降至 4.9540，10 轮优化计算约 58.9 秒，验证、采样和写盘单独计时。在相同训练预算下，42M、92M 的验证 NLL 进一步降至 4.8151 和 4.6831。扩容效果见图~\ref{fig:training}A。单卡短窗口测试的吞吐与显存见表~\ref{tab:capacity}，这些数据用于选择批量。

\begin{table}[!htbp]
\centering
\caption{固定唐诗语料与微批 32 的容量测试。显存为 PyTorch 峰值保留量。}
\label{tab:capacity}
\begin{tabular}{@{}lrrr@{}}
\toprule
模型 & 有效目标吞吐（token/s） & 峰值保留显存（GiB） & 10 轮验证 NLL\\
\midrule
14M & 456,778 & 1.88 & 4.9540\\
42M & 185,261 & 3.20 & 4.8151\\
92M & 101,098 & 4.97 & 4.6831\\
\bottomrule
\end{tabular}
\end{table}

预训练阶段已经学会生成古诗风格的短句。14M 的一次固定前缀续写为「{\versefont 春風自有情，獨坐故人歸。}」，随后生成结束符。内容和题意有一定联系，而输出只有两句。这项直接观察说明了预训练能提供字句基础，也给监督阶段提出了明确目标，即让输入词与完整四句建立稳定对应。

\subsection{延长训练与扩大唯一文本量}

只用唐诗的 92M 长程配置上限为 26 轮，实跑 16 轮、约 4,558 万有效目标。最佳验证 NLL 4.2995 出现在第 13 轮。此后三轮未改善并触发早停。复用相同的小语料能继续降低损失，但其收益会趋缓。与之配对分析的唐宋版扩大了唯一训练文本，并保持相同 92M 结构、字表与唐诗留出集。

\reportfigure{training.pdf}{预训练验证 NLL 随累计有效目标的变化。A 固定唐诗语料及目标数。B 比较唐诗长程与唐宋扩展，空心点标记各自最佳验证值。两条长程曲线的语料规模和学习率调度同时变化。}{fig:training}

唐宋版按 16 轮上限、每半轮验证训练。第 9 轮、第 5,076 步达到 NLL 3.8805，PPL 为 48.45，此时累计 166,066,074 个目标。训练于第 10.5 轮早停，处理 193,758,254 个目标，优化计算约 30.7 分钟。结合两条曲线可以看到，增加诗料和训练预算有益于同一留出集的语言建模。由于学习率计划和数据规模一起变化，这一结果只能归因于整套扩容配置。92M 微批从 32 提到 64 时，峰值保留显存升至 8.72 GiB，而短测吞吐从约 10.1 万降到约 9.8 万目标每秒。选择微批需要实测吞吐与稳定性。

\section{条件微调与韵律续训}
\label{sec:sft}

\subsection{监督标签与移位掩码}

公开语料缺少每首诗的人工关键词标签。先在训练集上统计意象词白名单，保留见于至少两首训练诗的词，排除停词、长词内包含的短词及跨句匹配。在每首合格诗中按固定种子选择一至三个词并打乱顺序，训练目标仍是原作的完整诗文。唐宋微调训练共有 12,394 对，纳入监督前剔除 776 条目标正文含未登录字符的匹配诗。这种弱监督能在一般的古诗中构造大量配对，但人类读题时常用的反衬、暗示和同义词不会直接进入字面标签。

输入前缀依次提供 \code{<bos>}、\code{<wujue>}、\code{<kw>}、关键词及 \code{<sep>}。训练时，前缀及填充的目标位置记为 \code{-100}。由于标签向前移一位，分隔符所在位置预测第一字，所以从该位置开始计损失，最后一个标点预测结束符。若把该位置也屏蔽，模型就失去了「看到条件后选首字」的直接反馈。图~\ref{fig:sftmask} 对应一条真实监督样本。

\reportfigure{supervision.pdf}{监督配对的输入和下一标记标签。长条件与长正文略去一部分，索引与目标来自冻结的掩码样例。}{fig:sftmask}

\subsection{从基线调通微调到唐宋版选模}

14M 从第 870 步的最佳预训练权重开始监督训练。开发阶段先比较峰值学习率 $3\times10^{-5}$ 与 $3\times10^{-4}$。在相同验证子集上，后者的末轮句式和全词联合达标由 3/30 升到 13/30，而损失仅有轻微反向变化。正式三种子训练共用预训练初始化，每个种子 5 轮、460 步及 366,125 个有效目标。种子 2027 为该阶段的演示权重，以便在保留的 14M 对照和盲评中追溯。

唐宋 92M 从第 9 轮最佳预训练权重开始，仍用三个种子分别创建优化器。扩大配对后采用 $1.5\times10^{-4}$ 的峰值学习率和最多 25 轮的调度，每个种子因验证早停实际运行 4 轮、1,552 步、1,239,400 个目标。三个最低监督验证 NLL 均出现在第 1 轮，分别为 3.6487、3.6600 和 3.6871，微调前的同口径值为 5.5932。先以各自最低 NLL 定种子内检查点，再按验证诗题的字数、用词及训练重合条件选择交付种子。90 道验证题每题采样三次，三个种子无整首或长诗片段命中的联合成功数分别为 66、79、95，因此选出 2028。选择方案在接触冻结测试结果前完成。

\subsection{同韵样本的匹配续训}

为改善自由生成中的押韵，从 92M 同一个最佳监督权重开始，对照抽取 4,096 条普通训练配对，处理组抽取 4,096 条偶句异字、同一平水韵平声部的配对~\cite{pingshui}。两组在词数上匹配，一词、两词、三词样本量分别为 2,028、1,342 和 726。两组用同一个训练种子、相同优化参数，各训练 6 轮、384 步、614,400 个目标。普通组和同韵组的数据取自原训练集合。在 90 道验证题各采 3 次的自由生成中，原权重、普通组、同韵组的异字同韵数分别为 117、119 和 188。

\reportfigure{finetune.pdf}{监督微调和韵律续训。A 为唐宋 SFT 三个种子的第 1 至第 4 轮验证 NLL。B 固定采样条件，三组的分母均为 270。}{fig:sft}

同韵组的 188/270 比原权重的 117/270 多 71 首。三项同时达标由 42/270 到 80/270，训练集整首匹配仍为 0，因此按事先约定切换到同韵组末权重。这是一组同训练种子的配对实验，对稳定性的估计有限。它支持筛选可靠韵脚样本可以改善当前模型的自由押韵，推广到其他语料或韵部仍需更多重复。
